\documentclass[12pt]{article}
\usepackage[ukrainian]{babel}
\usepackage{fontspec}
% --- НАЛАШТУВАННЯ ШРИФТІВ ---
\setmainfont{Schoolbook-Regular.otf}[
    Path = ../,
    BoldFont = Schoolbook-Bold.otf,
    ItalicFont = Schoolbook-Italic.otf,
    BoldItalicFont = Schoolbook-BoldItalic.otf
]
\usepackage[italic]{mathastext}
\usepackage[a4paper,margin=1.8cm,bottom=2cm,top=2cm]{geometry}
\usepackage{amsmath,amssymb}
\usepackage{tikz}
\usepackage{pgfplots}
\pgfplotsset{compat=1.16}
\usetikzlibrary{calc,patterns,angles,quotes,intersections}
\usepackage{xcolor,array,fancyhdr,enumitem,multicol}
\usepackage{colortbl}
\usepackage{tcolorbox}
\tcbuselibrary{skins,breakable}
\usepackage[hidelinks]{hyperref}
\usepackage[firstpage=false, text=@pvtr2525, scale=0.6, color=gray!12]{draftwatermark}

\binoppenalty=10000
\relpenalty=10000

\definecolor{mainGreen}{RGB}{34, 120, 64}
\definecolor{yearOrange}{RGB}{220, 100, 30}
\definecolor{titleBg}{RGB}{225, 240, 220}
\definecolor{instrBg}{RGB}{255, 240, 220}
\definecolor{instrBorder}{RGB}{220, 150, 80}
\definecolor{hexcol}{RGB}{200,110,60}

\pagestyle{fancy}
\fancyhf{}
\renewcommand{\headrulewidth}{0.4pt}
\renewcommand{\footrulewidth}{0.4pt}
\fancyhead[C]{\small\color{gray!80} Збірник НМТ \textendash{} Варіант за 17 червня}
\fancyhead[R]{\small\color{mainGreen}\textbf{\href{https://t.me/pvtr2525}{@pvtr2525}}}
\fancyfoot[L]{\small\color{gray!80}\href{https://t.me/pvtr2525}{ТГ: @pvtr2525}}
\fancyfoot[C]{\small\color{gray!80}\thepage}
\fancyfoot[R]{\small\color{gray!80}НМТ 2026}

\newcommand{\answerTable}[5]{%
\par\vspace{0.15cm}
\noindent
\begin{tabular}{|*{5}{>{\centering\arraybackslash}m{3cm}|}}
\hline
\rule[-0.15cm]{0pt}{0.6cm}\textbf{А} & \rule[-0.15cm]{0pt}{0.6cm}\textbf{Б} & \rule[-0.15cm]{0pt}{0.6cm}\textbf{В} & \rule[-0.15cm]{0pt}{0.6cm}\textbf{Г} & \rule[-0.15cm]{0pt}{0.6cm}\textbf{Д} \\
\hline
\rule[-0.2cm]{0pt}{0.75cm}#1 & \rule[-0.2cm]{0pt}{0.75cm}#2 & \rule[-0.2cm]{0pt}{0.75cm}#3 & \rule[-0.2cm]{0pt}{0.75cm}#4 & \rule[-0.2cm]{0pt}{0.75cm}#5 \\
\hline
\end{tabular}
\par\vspace{0.25cm}
}

\newcommand{\answerTableTall}[5]{%
\par\vspace{0.15cm}
\noindent
\begin{tabular}{|*{5}{>{\centering\arraybackslash}m{3cm}|}}
\hline
\rule[-0.2cm]{0pt}{0.7cm}\textbf{А} & \rule[-0.2cm]{0pt}{0.7cm}\textbf{Б} & \rule[-0.2cm]{0pt}{0.7cm}\textbf{В} & \rule[-0.2cm]{0pt}{0.7cm}\textbf{Г} & \rule[-0.2cm]{0pt}{0.7cm}\textbf{Д} \\
\hline
\rule[-0.45cm]{0pt}{1.2cm}#1 & \rule[-0.45cm]{0pt}{1.2cm}#2 & \rule[-0.45cm]{0pt}{1.2cm}#3 & \rule[-0.45cm]{0pt}{1.2cm}#4 & \rule[-0.45cm]{0pt}{1.2cm}#5 \\
\hline
\end{tabular}
\par\vspace{0.25cm}
}

\newcommand{\instructionBox}[1]{%
\par\vspace{0.3cm}
\begin{tcolorbox}[colback=instrBg, colframe=instrBorder, boxrule=0.6pt, arc=2pt,
    left=8pt, right=8pt, top=4pt, bottom=4pt]
\centering\small\bfseries #1
\end{tcolorbox}
\par\vspace{0.15cm}
}

\newcommand{\sectionTitle}[1]{%
\par\vspace{0.4cm}
\begin{tcolorbox}[colback=titleBg, colframe=titleBg, boxrule=0pt, arc=8pt, halign=center, fontupper=\Large\bfseries\color{mainGreen}]
#1
\end{tcolorbox}\par\vspace{0.3cm}}
\newcommand{\chapterTitle}[1]{\clearpage\sectionTitle{#1}}

\newcommand{\nmtAnswerBox}{%
\par\vspace{0.2cm}\noindent
Відповідь:\ \ 
\begin{tabular}{@{}*{4}{|>{\centering\arraybackslash}p{0.8cm}}|@{\hspace{0.1cm}\textbf{,}\hspace{0.1cm}}*{3}{|>{\centering\arraybackslash}p{0.8cm}}|@{}}
\hline
\rule[-0.2cm]{0pt}{0.7cm} & & & & & & \\
\hline
\end{tabular}
\par\vspace{0.3cm}}

\newcommand{\matchingGrid}{%
    \begingroup
    \renewcommand{\arraystretch}{1.15}
    \setlength{\tabcolsep}{4pt}
    \footnotesize
    \begin{tabular}{r|c|c|c|c|c|}
         \multicolumn{1}{c}{} & \multicolumn{1}{c}{\textbf{А}} & \multicolumn{1}{c}{\textbf{Б}} & \multicolumn{1}{c}{\textbf{В}} & \multicolumn{1}{c}{\textbf{Г}} & \multicolumn{1}{c}{\textbf{Д}} \\ \cline{2-6}
         \textbf{1} & \phantom{X} & \phantom{X} & \phantom{X} & \phantom{X} & \phantom{X} \\ \cline{2-6}
         \textbf{2} & & & & & \\ \cline{2-6}
         \textbf{3} & & & & & \\ \cline{2-6}
    \end{tabular}
    \endgroup
}

% flat-top шестикутник зі стороною 30 (для №3)
\newcommand{\hexthree}[2]{%
    \draw[hexcol, line width=1.1pt]
        ($(#1,#2)+(0:30)$) -- ($(#1,#2)+(60:30)$) -- ($(#1,#2)+(120:30)$) --
        ($(#1,#2)+(180:30)$) -- ($(#1,#2)+(240:30)$) -- ($(#1,#2)+(300:30)$) -- cycle;
}

\newcounter{zad}
\newcommand{\zadtask}[1]{\stepcounter{zad}\par\vspace{0.3cm}\noindent\textbf{\thezad.}\ \ #1\par\vspace{0.2cm}}
\newcommand{\zadnum}{\stepcounter{zad}\noindent\textbf{\thezad.}\ \ }

\begin{document}
\thispagestyle{empty}
\vspace*{1.2cm}
\begin{center}
{\Large\bfseries\color{mainGreen} РЕАЛЬНИЙ ВАРІАНТ НМТ-2026}\\[0.4cm]
{\Huge\bfseries\color{mainGreen} ВАРІАНТ ЗА 17 ЧЕРВНЯ}\\[0.6cm]
{\large Real Test Notebook з усіма геометричними рисунками}\\[1cm]
{\large\color{mainGreen}\textbf{\href{https://t.me/pvtr2525}{Телеграм: @pvtr2525}}}
\end{center}
\clearpage

\chapterTitle{ТЕСТОВИЙ ЗОШИТ (17 ЧЕРВНЯ)}
\instructionBox{Завдання 1--15 мають по п'ять варіантів відповіді, з яких лише один правильний. Виберіть правильний варіант відповіді й позначте його.}

% ===== №1 =====
\zadtask{У магазині сухофруктів сушений лісовий горіх коштує $420$ гривень за один кілограм. Андрій купив $300$ грамів таких горіхів. Скільки грошей заплатив Андрій за покупку?}
\answerTable{$120$~грн}{$126$~грн}{$12{,}6$~грн}{$140$~грн}{$14$~грн}

% ===== №2 =====
\noindent
\begin{minipage}[c]{0.46\textwidth}
\zadnum На діаграмі відображено частку відвідувачів (у~\%), які позитивно оцінили якість обслуговування в супермаркеті протягом п'яти місяців року. Визначте місяць, у якому ця частка становила $\dfrac{4}{5}$ від загальної кількості опитаних відвідувачів у цьому місяці.
\end{minipage}%
\hfill
\begin{minipage}[c]{0.52\textwidth}
\centering
\begin{tikzpicture}[x=0.9cm, y=0.32cm, thick]
    \draw[->] (0,0) -- (0,11) node[above, font=\scriptsize, rotate=90, anchor=south, xshift=-2.2cm, yshift=0.4cm] {Частка відвідувачів, \%};
    \draw[->] (0,0) -- (5.6,0);
    \foreach \y in {1,2,3,4,5,6,7,8,9,10} {
        \draw[gray!25, thin] (0,\y) -- (5.3,\y);
        \node[left, font=\scriptsize] at (0,\y) {\pgfmathparse{int(\y*10)}\pgfmathresult};
    }
    \node[left, font=\scriptsize] at (0,0) {0};
    \definecolor{cjan}{RGB}{205,205,215}
    \definecolor{cfeb}{RGB}{95,190,225}
    \definecolor{cmar}{RGB}{120,175,55}
    \definecolor{capr}{RGB}{150,50,120}
    \definecolor{cmay}{RGB}{220,45,50}
    \filldraw[fill=cjan, draw=none] (0.55,0) rectangle ++(0.5,5);
    \filldraw[fill=cfeb, draw=none] (1.55,0) rectangle ++(0.5,8);
    \filldraw[fill=cmar, draw=none] (2.55,0) rectangle ++(0.5,6);
    \filldraw[fill=capr, draw=none] (3.55,0) rectangle ++(0.5,9);
    \filldraw[fill=cmay, draw=none] (4.55,0) rectangle ++(0.5,7);
    \node[below, font=\tiny] at (0.8,0) {Січень};
    \node[below, font=\tiny] at (1.8,0) {Лютий};
    \node[below, font=\tiny] at (2.8,0) {Березень};
    \node[below, font=\tiny] at (3.8,0) {Квітень};
    \node[below, font=\tiny] at (4.8,0) {Травень};
\end{tikzpicture}
\end{minipage}
\par\vspace{0.5cm}
\answerTable{січень}{лютий}{березень}{квітень}{травень}

% ===== №3 =====
\noindent
\begin{minipage}[c]{0.52\textwidth}
\zadnum На рисунку зображено настінну полицю, що складається з трьох однакових полиць у формі правильних шестикутників. Визначте загальну ширину $AB$ цієї конструкції, якщо сторона кожного шестикутника дорівнює $30$~см.
\end{minipage}%
\hfill
\begin{minipage}[c]{0.46\textwidth}
\centering
\begin{tikzpicture}[scale=0.05, line join=round]
    \hexthree{0}{0}
    \hexthree{90}{0}
    \hexthree{45}{-25.98}
    \coordinate (A) at (-30,0);
    \coordinate (B) at (120,0);
    \node[left, font=\small] at (A) {$A$};
    \node[right, font=\small] at (B) {$B$};
    \draw[<->] (-30,-54) -- (120,-54);
    \draw[thin, gray] (-30,0) -- (-30,-54);
    \draw[thin, gray] (120,0) -- (120,-54);
    \node[below, font=\small] at (45,-54) {?};
\end{tikzpicture}
\end{minipage}
\par\vspace{0.2cm}
\answerTable{$60$~см}{$90$~см}{$120$~см}{$150$~см}{$180$~см}

% ===== №4 =====
\zadtask{Укажіть пару чисел, що є розв'язком системи рівнянь $\begin{cases}xy = 32,\\ y = \dfrac{x}{2}.\end{cases}$}
\answerTableTall{$(4; 8)$}{$(16; 2)$}{$(-4; -8)$}{$(2; 4)$}{$(8; 4)$}

% ===== №5 =====
\zadtask{У якій просторовій фігурі можна провести твірну?
\vspace{0.1cm}
\begin{enumerate}[label=\textbf{\Alph*}, leftmargin=2.8em, itemsep=0.1cm]
    \item лише в конусі
    \item лише в піраміді
    \item лише в конусі та піраміді
    \item лише в конусі та циліндрі
    \item лише в циліндрі
\end{enumerate}
\vspace{0.1cm}}

% ===== №6 =====
\zadtask{$\dfrac{3a - b}{18a^2 - 2b^2} = $}
\answerTableTall{$6a + 2b$}{$6a - 2b$}{$\dfrac{1}{6a + 2b}$}{$\dfrac{1}{2}$}{$\dfrac{1}{6a - 2b}$}

% ===== №7 =====
\zadtask{Розв'яжіть рівняння $3^x = 9^{x+2}$.}
\answerTable{$-4$}{$4$}{$0$}{$-2$}{$2$}

% ===== №8 =====
\noindent
\begin{minipage}[c]{0.58\textwidth}
\zadnum На рисунку зображено графік квадратичної функції. Укажіть точку локального екстремуму цієї функції.
\end{minipage}%
\hfill
\begin{minipage}[c]{0.40\textwidth}
\centering
\begin{tikzpicture}[scale=0.5, thick]
    \draw[step=1cm, gray!40, thin] (-2,-3) grid (4,4);
    \draw[->] (-2.3,0) -- (4.3,0) node[right, font=\scriptsize] {$x$};
    \draw[->] (0,-3.3) -- (0,4.3) node[above, font=\scriptsize] {$y$};
    \node[below left, font=\scriptsize] at (0,0) {$0$};
    \node[below, font=\scriptsize] at (1,0) {$1$};
    \node[left, font=\scriptsize] at (0,1) {$1$};
    % парабола гілками вниз, вершина (1,3): y = -(x-1)^2+3 -> прибл
    \draw[mainGreen, very thick, domain=-1.05:3.05, samples=60] plot (\x, {-1*(\x-1)*(\x-1)+3});
\end{tikzpicture}
\end{minipage}
\par\vspace{0.2cm}
\answerTable{$x_0 = 4$}{$x_0 = 0$}{$x_0 = 1$}{$x_0 = 3$}{$x_0 = -1$}

% ===== №9 =====
\zadtask{Які з наведених тверджень щодо кола є правильними?
\vspace{0.1cm}
\begin{enumerate}[label=\textbf{\Roman*.}, align=left, leftmargin=2.6em, labelsep=0.5em, itemsep=0.08cm, topsep=0.06cm]
\item Дуга кола~--- це відрізок, який сполучає $2$ точки на колі.
\item Центральний кут~--- це кут, вершина якого лежить на колі.
\item Вписаний кут, що спирається на діаметр,~--- прямий.
\end{enumerate}
\vspace{0.1cm}}
\answerTable{лише I}{лише II}{лише III}{лише I та III}{лише II та III}

% ===== №10 =====
\zadtask{У геометричній прогресії $(b_n)$ відомо, що $b_3 = 2b_2$, $b_2 \neq 0$. Знайдіть значення виразу $\dfrac{b_5}{b_2}$.}
\answerTableTall{$\dfrac{1}{6}$}{$4$}{$6$}{$8$}{$16$}
\newpage

% ===== №11 =====
\zadtask{Дмитро хоче придбати на офіційному сайті «Укрпошти» один екземпляр колекційної марки. У наявності інтернет-магазину є $10$ марок серії «Міста-герої», $20$ марок серії «Славетна Україна» та $14$ марок серії «Культурна спадщина». Скільки всього є варіантів вибору однієї марки у Дмитра, якщо всі наявні марки є різними?}
\answerTable{$34$}{$44$}{$280$}{$2400$}{$2800$}

% ===== №12 =====
\noindent
\begin{minipage}[c]{0.56\textwidth}
\zadnum У прямокутній системі координат у просторі зображено куб $ABCDA_1B_1C_1D_1$, вершина $B$ якого збігається з початком координат, а ребра $AB$, $BC$ й $BB_1$ лежать на координатних осях (див. рисунок). Ребро куба дорівнює $1$. Визначте координати вектора $\overrightarrow{AD_1}$.
\vspace{0.2cm}
\par\noindent\textbf{А}\ \ $(1; 1; 1)$
\par\noindent\textbf{Б}\ \ $(1; 1; 0)$
\par\noindent\textbf{В}\ \ $(1; 0; 1)$
\par\noindent\textbf{Г}\ \ $(1; 0; 0)$
\par\noindent\textbf{Д}\ \ $(0; 1; 1)$
\end{minipage}%
\hfill
\begin{minipage}[c]{0.42\textwidth}
\centering
\begin{tikzpicture}[scale=1.5, line join=round]
    \coordinate (B) at (0,0);
    \coordinate (C) at (1.3,0);
    \coordinate (A) at (-0.55,-0.55);
    \coordinate (D) at (0.75,-0.55);
    \coordinate (B1) at (0,1.3);
    \coordinate (C1) at (1.3,1.3);
    \coordinate (A1) at (-0.55,0.75);
    \coordinate (D1) at (0.75,0.75);
    % осі
    \draw[->] (0,0) -- (-1.0,-1.0) node[below, font=\scriptsize] {$x$};
    \draw[->] (0,0) -- (2.0,0) node[right, font=\scriptsize] {$y$};
    \draw[->] (0,0) -- (0,2.0) node[above, font=\scriptsize] {$z$};
    % ребра куба
    \draw[dashed] (A) -- (B) -- (C);
    \draw[dashed] (B) -- (B1);
    \draw (A) -- (D) -- (C);
    \draw (A) -- (A1);
    \draw (D) -- (D1);
    \draw (C) -- (C1);
    \draw (A1) -- (B1) -- (C1) -- (D1) -- cycle;
    \node[below left, font=\scriptsize] at (A) {$A$};
    \node[below left, font=\scriptsize] at (B) {$B$};
    \node[right, font=\scriptsize] at (C) {$C$};
    \node[below right, font=\scriptsize] at (D) {$D$};
    \node[left, font=\scriptsize] at (A1) {$A_1$};
    \node[above left, font=\scriptsize] at (B1) {$B_1$};
    \node[above right, font=\scriptsize] at (C1) {$C_1$};
    \node[right, font=\scriptsize] at (D1) {$D_1$};
\end{tikzpicture}
\end{minipage}
\par\vspace{0.3cm}

% ===== №13 =====
\zadtask{Спростіть вираз $\dfrac{1 - \cos^2 \alpha}{\operatorname{tg} \alpha \cdot \cos \alpha}$.}
\answerTable{$\sin \alpha$}{$\operatorname{tg} \alpha$}{$1$}{$\cos \alpha$}{$0$}

% ===== №14 =====
\noindent
\begin{minipage}[c]{0.56\textwidth}
\zadnum У паралелограмі $ABCD$ з вершини тупого кута $B$ проведено висоти $BK$ і $BM$, кут між якими дорівнює $60^\circ$ (див. рисунок). $BC = 12$~см, $BK = 4\sqrt{3}$~см. Визначте довжину діагоналі $BD$ цього паралелограма.
\end{minipage}%
\hfill
\begin{minipage}[c]{0.42\textwidth}
\centering
\begin{tikzpicture}[scale=0.72, thick, line join=round]
    \coordinate (A) at (0,0);
    \coordinate (D) at (4,0);
    \coordinate (B) at (1.4,2.6);
    \coordinate (C) at (5.4,2.6);
    \coordinate (K) at (1.4,0);
    % M на CD, підніжжя перпендикуляра з B
    \coordinate (M) at ($(D)!(B)!(C)$);
    \draw (A) -- (B) -- (C) -- (D) -- cycle;
    \draw (B) -- (K);
    \draw (B) -- (M);
    \pic[draw, angle radius=0.18cm] {right angle = B--K--A};
    \pic[draw, angle radius=0.18cm] {right angle = D--M--B};
    \pic[draw, angle radius=0.4cm, pic text={$60^\circ$}, angle eccentricity=1.5, font=\scriptsize] {angle = K--B--M};
    \node[below left] at (A) {$A$};
    \node[above left] at (B) {$B$};
    \node[above right] at (C) {$C$};
    \node[below right] at (D) {$D$};
    \node[below] at (K) {$K$};
    \node[right] at (M) {$M$};
\end{tikzpicture}
\end{minipage}
\par\vspace{0.2cm}
\answerTable{$16$~см}{$4\sqrt{7}$~см}{$4$~см}{$8$~см}{$8\sqrt{3}$~см}

% ===== №15 =====
\zadtask{Знайдіть найменший цілий розв'язок нерівності $(\log_3 1 - 2)(3x - 6) \leqslant -6$.}
\answerTable{$5$}{$4$}{$3$}{$2$}{$1$}
\newpage
\instructionBox{У завданнях 16--18 до кожного з трьох рядків інформації, позначених цифрами, доберіть один правильний, на Вашу думку, варіант, позначений буквою.}

% ===== №16 =====
\zadtask{Узгодьте вираз (1--3) з його значенням (А--Д).\par
\vspace{0.3cm}
\noindent
\begin{minipage}[t]{0.42\textwidth}
\textit{Вираз}\par\vspace{0.15cm}
\textbf{1} \quad $|-2| \cdot 2^0$\\[0.35cm]
\textbf{2} \quad $\sqrt[3]{-\dfrac{1}{8}}$\\[0.35cm]
\textbf{3} \quad $3^{-\log_3 2}$
\end{minipage}%
\hfill
\begin{minipage}[t]{0.3\textwidth}
\textit{Значення виразу}\par\vspace{0.15cm}
\textbf{А} \quad $-2$\\[0.15cm]
\textbf{Б} \quad $-0{,}5$\\[0.15cm]
\textbf{В} \quad $0$\\[0.15cm]
\textbf{Г} \quad $0{,}5$\\[0.15cm]
\textbf{Д} \quad $2$
\end{minipage}%
\hfill
\begin{minipage}[t]{0.2\textwidth}
\vspace{0.4cm}
\begin{flushright}\matchingGrid\end{flushright}
\end{minipage}}

% ===== №17 =====
\zadtask{Доберіть до функції (1--3) властивість її графіка (А--Д).\par
\vspace{0.3cm}
\noindent
\begin{minipage}[t]{0.32\textwidth}
\textit{Функція}\par\vspace{0.2cm}
\textbf{1} \quad $y = 2 - x$\\[0.35cm]
\textbf{2} \quad $y = \dfrac{1}{x} - 2$\\[0.35cm]
\textbf{3} \quad $y = \log_{0{,}5} x$
\par\vspace{0.4cm}
\matchingGrid
\end{minipage}%
\hfill
\begin{minipage}[t]{0.64\textwidth}
\textit{Властивість графіка функції}\par\vspace{0.2cm}
\textbf{А} \quad перетинає вісь $y$ з додатною ординатою, а вісь $x$ не перетинає\\[0.15cm]
\textbf{Б} \quad не перетинає ані вісь $x$, ані вісь $y$\\[0.15cm]
\textbf{В} \quad розміщений лише в I, II та IV координатних чвертях\\[0.15cm]
\textbf{Г} \quad перетинає вісь $y$ з від'ємною ординатою, а вісь $x$ не перетинає\\[0.15cm]
\textbf{Д} \quad перетинає вісь $x$, а вісь $y$ не перетинає
\end{minipage}}

% ===== №18 =====
\zadtask{Установіть відповідність між трикутником $ABC$ (1--3) та його площею (А--Д).
\vspace{0.3cm}

\noindent
\begin{minipage}[t]{0.56\textwidth}
\textit{Трикутник $ABC$}\par\vspace{0.35cm}
% 1
\noindent\textbf{1}\hspace{0.3cm}\begin{tikzpicture}[baseline=0cm, scale=0.55, thick, line join=round]
    \coordinate (A) at (0,0);
    \coordinate (C) at (4.6,0);
    \coordinate (B) at (2.9,0.9);
    \draw (A) -- (B) -- (C) -- cycle;
    \pic[draw, angle radius=0.12cm, pic text={\tiny $150^\circ$}, angle eccentricity=1.9] {angle = A--B--C};
    \node[below left, font=\scriptsize] at (A) {$A$};
    \node[above, font=\scriptsize] at (B) {$B$};
    \node[below right, font=\scriptsize] at (C) {$C$};
    \node[right, font=\scriptsize] at (5.0,0.5) {$AB\cdot BC = 48,\ \angle B = 150^\circ$};
\end{tikzpicture}
\par\vspace{0.5cm}
% 2
\noindent\textbf{2}\hspace{0.3cm}\begin{tikzpicture}[baseline=0cm, scale=0.5, thick, line join=round]
    \coordinate (A) at (0,0);
    \coordinate (C) at (5,0);
    \coordinate (H) at (2,0);
    \coordinate (B) at (2,2.4);
    \draw (A) -- (B) -- (C) -- cycle;
    \draw (B) -- (H);
    \pic[draw, angle radius=0.2cm] {right angle = B--H--C};
    \node[below left, font=\scriptsize] at (A) {$A$};
    \node[above, font=\scriptsize] at (B) {$B$};
    \node[below right, font=\scriptsize] at (C) {$C$};
    \node[below, font=\scriptsize] at (H) {$H$};
    \node[right, font=\scriptsize] at (5.3,1.2) {$BH = 4,\ AC = 9$};
\end{tikzpicture}
\par\vspace{0.5cm}
% 3
\noindent\textbf{3}\hspace{0.3cm}\begin{tikzpicture}[baseline=0cm, scale=0.55, thick, line join=round]
    \coordinate (A) at (0,0);
    \coordinate (C) at (3.2,0);
    \coordinate (B) at (0,3.2);
    \coordinate (O) at (0,1.6);
    \coordinate (K) at (1.6,1.6);
    \draw (A) -- (B) -- (C) -- cycle;
    \draw (O) -- (K);
    \pic[draw, angle radius=0.2cm] {right angle = O--A--C};
    \pic[draw, angle radius=0.2cm] {right angle = B--O--K};
    \path (B) -- (K) node[midway, sloped, font=\tiny] {$|$};
    \path (K) -- (C) node[midway, sloped, font=\tiny] {$|$};
    \node[below left, font=\scriptsize] at (A) {$A$};
    \node[above, font=\scriptsize] at (B) {$B$};
    \node[below right, font=\scriptsize] at (C) {$C$};
    \node[left, font=\scriptsize] at (O) {$O$};
    \node[above right, font=\scriptsize] at (K) {$K$};
    \node[right, font=\scriptsize] at (3.4,2.2) {$AO = 3,\ OK = 4$};
\end{tikzpicture}
\end{minipage}%
\hfill
\begin{minipage}[t]{0.24\textwidth}
\textit{Площа $ABC$}\par\vspace{0.2cm}
\textbf{А} \quad $12$\\[0.3cm]
\textbf{Б} \quad $18$\\[0.3cm]
\textbf{В} \quad $24$\\[0.3cm]
\textbf{Г} \quad $36$\\[0.3cm]
\textbf{Д} \quad $48$
\end{minipage}%
\hfill
\begin{minipage}[t]{0.16\textwidth}
\vspace{0.4cm}
\begin{flushright}\matchingGrid\end{flushright}
\end{minipage}
\par\vspace{0.2cm}}
\newpage
\instructionBox{Розв'яжіть завдання 19--22. Одержані числові відповіді запишіть у спеціально відведеному місці. Відповідь записуйте лише десятковим дробом, урахувавши положення коми. Знак «мінус» записуйте перед першою цифрою числа.}

% ===== №19 =====
\zadtask{Обчисліть $\displaystyle\int\limits_{1}^{4} \sqrt{x}\left(\sqrt{x} + 6\right)dx$.}
\nmtAnswerBox

% ===== №20 =====
\zadtask{Автомобіль двічі заправляли пальним і щоразу по $40$~л. Порівняно з першим заправленням ціна пального, використаного для другого заправлення, була більшою на $2{,}5\,\%$. Усього за ці два заправлення автомобіля пальним було витрачено $4860$~грн. Скільки \textit{гривень} коштував $1$~л пального, використаного під час першого заправлення?}
\nmtAnswerBox

% ===== №21 =====
\noindent
\begin{minipage}[c]{0.56\textwidth}
\zadnum На рисунку зображено правильну трикутну призму й конус. Твірна конуса дорівнює $12$~см, а площа його бічної поверхні дорівнює $48\pi$~см$^2$. Радіус основи конуса дорівнює радіусу кола, описаного навколо основи призми. Знайдіть об'єм (у см$^3$) призми, якщо її бічні грані є квадратами.
\end{minipage}%
\hfill
\begin{minipage}[c]{0.42\textwidth}
\centering
\begin{tikzpicture}[scale=0.6, thick, line join=round]
    % трикутна призма
    \coordinate (a) at (0,0);
    \coordinate (b) at (2.4,0.5);
    \coordinate (c) at (1.3,1.4);
    \coordinate (a2) at (0,4);
    \coordinate (b2) at (2.4,4.5);
    \coordinate (c2) at (1.3,5.4);
    \draw (a) -- (b) -- (c) -- cycle;
    \draw (a2) -- (b2) -- (c2) -- cycle;
    \draw (a) -- (a2);
    \draw (b) -- (b2);
    \draw[dashed] (c) -- (c2);
    % конус
    \begin{scope}[shift={(4.4,0)}]
        \draw[dashed] (-1.3,0.5) arc (180:0:1.3 and 0.4);
        \draw (-1.3,0.5) arc (180:360:1.3 and 0.4);
        \draw (-1.3,0.5) -- (0,4.6);
        \draw (1.3,0.5) -- (0,4.6);
    \end{scope}
\end{tikzpicture}
\end{minipage}
\par\vspace{0.2cm}
\nmtAnswerBox

% ===== №22 =====
\zadtask{Знайдіть \textit{суму} всіх цілих значень $a$, за кожного з яких рівняння
\[
\dfrac{\log_3(x^2 + 3x - 1) - 2}{\sqrt{x + a} + \sqrt{10 - a}} = 0
\]
має два різні корені.}
\nmtAnswerBox

\newpage
% ===== ВІДПОВІДІ =====

% ===== ВІДПОВІДІ (офіційні) =====
\sectionTitle{ВІДПОВІДІ ДО ВАРІАНТА (17 ЧЕРВНЯ)}
\textbf{1.}\ Б\\
\textbf{2.}\ Б\\
\textbf{3.}\ Г\\
\textbf{4.}\ Д\\
\textbf{5.}\ Г\\
\textbf{6.}\ В\\
\textbf{7.}\ А\\
\textbf{8.}\ В\\
\textbf{9.}\ В\\
\textbf{10.}\ Г\\
\textbf{11.}\ Б\\
\textbf{12.}\ Д\\
\textbf{13.}\ А\\
\textbf{14.}\ Б\\
\textbf{15.}\ Б\\
\textbf{16.}\ 1--Д, 2--Б, 3--Г\\
\textbf{17.}\ 1--В, 2--Г, 3--Д\\
\textbf{18.}\ 1--А, 2--Б, 3--В\\
\textbf{19.}\ $35.5$\\
\textbf{20.}\ $60$\\
\textbf{21.}\ $144$\\
\textbf{22.}\ $45$\\

\end{document}
